%-----------------------------------------------------------------------
\section{Limitations}
\label{sec:limitations}
%-----------------------------------------------------------------------

\subsection{The correction is defined for regular fold arcs}
\label{sec:limits-fold}

The profile of Section~\ref{sec:ours-normalisation} is the diffraction pattern
of a fold ($A_2$). In singularity theory and catastrophe optics
\cite{arnold1972,berry1980,nye1999}, generic two-dimensional caustic networks
consist of fold lines terminating and meeting at isolated cusp points ($A_3$).
Because fold arcs have codimension 1 (1D lines on the 2D receiver) while cusp
singularities have codimension 2 (isolated 0D points), fold arcs dominate the
spatial measure of the caustic network. Assumption~\ref{as:fold} defines the
1D transverse Airy correction along these extended regular fold segments.

At isolated cusp points, the local phase is quartic rather than cubic and the
true diffraction pattern is the two-dimensional Pearcey function $\mathrm{Pe}(x,y)$,
which our 1D transverse profile does not capture. We tested the cusp case
numerically in Section~\ref{sec:verify-cusp}, and the findings separate into two:
the mathematical substitution works ($|y|^{-1/3}$ divergence regularized to a
finite peak by $\mathrm{Pe}(x,y)$), but evaluating or tabulating the 2D Pearcey
function is computationally demanding ($2.1 \times 10^4$ slower than a 1D Airy
table lookup). Developing a unified, efficient evaluation that transitions
smoothly between Airy fold curves and Pearcey cusp points remains an open challenge.
Swallowtails and umbilic points remain untested.

\subsection{The coherent limit is not the sun}
\label{sec:limits-coherence}

This is the most serious physical limitation. The source in this scene subtends
\SI{0.53}{\degree}, and at the receiver distance that corresponds to a geometric
blur of roughly \SI{925}{\micro\meter}. The physical fringe spacing implied by
\eqref{eq:scalelaw} for this geometry is \SI{123}{\micro\meter}, about
\num{7.5} times smaller. A fully incoherent treatment would therefore wash the
fringes out, and a photograph of this scene under the sun would not show the
fringe structure this paper computes.

What the paper's fringe results therefore establish is that the wave-corrected
estimator reproduces the \emph{coherent-limit} fold profile, that the profile is
the Airy function, and that the resulting field is bounded and smooth where the
geometric estimator is neither. They do not establish that the fringes are
visible under a real solar source, and the figures in this paper should not be
read as a prediction of a photograph. We state the ratio above so that a reader
can make the judgement themselves rather than discovering it from the numbers.

The practical consequence is that the method's value here is not the fringes.
It is that replacing a divergent integrand with a bounded one, at the correct
flux, produces a caustic that is both closer to the converged answer and stable
in time (Section~\ref{sec:results-flicker}). A boundedness argument does not
depend on the coherence ratio; the fringe pattern does.

\subsection{Not resolved at the production grid}
\label{sec:limits-resolution}

The fold's apparent breadth measures $2.30\times$ the geometric pixel width at
$1024^2$ and $1.87\times$ at $4096^2$. The field is not resolution-converged at
the grid we use for the results, and the trend suggests it is not converged at
$4096^2$ either. Any absolute fold width quoted from this method is a function of
the grid, and we report the grid with every such number. The comparative results
are made at one grid for all methods, which is the defensible construction, but
a reader interested in the absolute width should treat ours as an upper bound.

\subsection{The ceiling is a safeguard}
\label{sec:limits-ceiling}

The display ceiling $\kappa$ is an optional safeguard against runaway geometric
outlier splats. As shown in the sensitivity study (Table~\ref{tab:sensitivity}),
the wave correction naturally bounds the caustic peak to $\approx 3.1\times$ the
mean, so setting $\kappa \ge 4.0\times$ leaves the regularized fold flux completely
untouched ($\text{RMSE} = 0.0000$). Band flux is strictly conserved by
Lemma~\ref{lem:flux} provided $\kappa$ is not set below the physical Airy
peak.

\subsection{Estimator dependence}
\label{sec:limits-estimators}

Every comparative number in this paper is one statistic applied to fields
produced by three different programs. We chose the statistics before the
comparison (Section~\ref{sec:estimators}), which is the correct order, and we
report the code that produces them so the choice can be reversed. But a claim
that flips under a change of statistic is not a result, and we found exactly one
such statistic in our own history
(Section~\ref{sec:verify-metric}). Others may remain, and the reader is
entitled to check.

\subsection{No physical apparatus}
\label{sec:limits-lab}

The comparison is entirely computational. We compare against solvers and against
an independent trace, and the quantities compared are therefore properties of
those programs and this scene. A laboratory test would address a different
claim --- whether the geometric-plus-Airy model predicts a measured intensity
distribution --- and it would additionally test the coherence question of
Section~\ref{sec:limits-coherence} in the only way that question can be settled.
The scene is specified precisely enough in Table~\ref{tab:scene} that this is
feasible, and we regard it as the natural continuation.

\subsection{Scope of the cost claims}
\label{sec:limits-cost}

The timing figures are from one machine and one implementation. The method's
advantage is a boundedness argument with a cost attached, and the cost scales
with the band's measure, which is scene-dependent. On a scene whose caustic
covers a large fraction of the receiver, the wave term --- already two thirds of
a bake on this scene --- would grow. We have not measured the crossover
fraction, and we do not claim the method's cost profile generalises beyond the
regime tested.

\subsection{What would falsify the approach}
\label{sec:limits-falsify}

To keep the claims checkable we state what we would take as a refutation. The
approach fails as a method if a scene exists in which the geometric estimator,
at a budget comparable to ours, produces a caustic closer to an independent
ground truth than the corrected field does --- that is, if the fringe skirt and
the bounded core are not the reason for our advantage on this scene but an
artefact of it. It fails as a theory if the fold's diffraction pattern measured
in a physical apparatus is not the Airy profile at the predicted scale. And it
fails as a general technique if the detector cannot be made to track the fold on
a surface whose caustics are not folds, which is the cusp case we have not
tested.
